% Options for packages loaded elsewhere
\PassOptionsToPackage{unicode}{hyperref}
\PassOptionsToPackage{hyphens}{url}
%
\documentclass[
]{book}
\usepackage{lmodern}
\usepackage{amssymb,amsmath}
\usepackage{ifxetex,ifluatex}
\ifnum 0\ifxetex 1\fi\ifluatex 1\fi=0 % if pdftex
  \usepackage[T1]{fontenc}
  \usepackage[utf8]{inputenc}
  \usepackage{textcomp} % provide euro and other symbols
\else % if luatex or xetex
  \usepackage{unicode-math}
  \defaultfontfeatures{Scale=MatchLowercase}
  \defaultfontfeatures[\rmfamily]{Ligatures=TeX,Scale=1}
\fi
% Use upquote if available, for straight quotes in verbatim environments
\IfFileExists{upquote.sty}{\usepackage{upquote}}{}
\IfFileExists{microtype.sty}{% use microtype if available
  \usepackage[]{microtype}
  \UseMicrotypeSet[protrusion]{basicmath} % disable protrusion for tt fonts
}{}
\makeatletter
\@ifundefined{KOMAClassName}{% if non-KOMA class
  \IfFileExists{parskip.sty}{%
    \usepackage{parskip}
  }{% else
    \setlength{\parindent}{0pt}
    \setlength{\parskip}{6pt plus 2pt minus 1pt}}
}{% if KOMA class
  \KOMAoptions{parskip=half}}
\makeatother
\usepackage{xcolor}
\IfFileExists{xurl.sty}{\usepackage{xurl}}{} % add URL line breaks if available
\IfFileExists{bookmark.sty}{\usepackage{bookmark}}{\usepackage{hyperref}}
\hypersetup{
  pdftitle={Marketing Research},
  pdfauthor={Mike Nguyen},
  hidelinks,
  pdfcreator={LaTeX via pandoc}}
\urlstyle{same} % disable monospaced font for URLs
\usepackage{color}
\usepackage{fancyvrb}
\newcommand{\VerbBar}{|}
\newcommand{\VERB}{\Verb[commandchars=\\\{\}]}
\DefineVerbatimEnvironment{Highlighting}{Verbatim}{commandchars=\\\{\}}
% Add ',fontsize=\small' for more characters per line
\usepackage{framed}
\definecolor{shadecolor}{RGB}{248,248,248}
\newenvironment{Shaded}{\begin{snugshade}}{\end{snugshade}}
\newcommand{\AlertTok}[1]{\textcolor[rgb]{0.94,0.16,0.16}{#1}}
\newcommand{\AnnotationTok}[1]{\textcolor[rgb]{0.56,0.35,0.01}{\textbf{\textit{#1}}}}
\newcommand{\AttributeTok}[1]{\textcolor[rgb]{0.77,0.63,0.00}{#1}}
\newcommand{\BaseNTok}[1]{\textcolor[rgb]{0.00,0.00,0.81}{#1}}
\newcommand{\BuiltInTok}[1]{#1}
\newcommand{\CharTok}[1]{\textcolor[rgb]{0.31,0.60,0.02}{#1}}
\newcommand{\CommentTok}[1]{\textcolor[rgb]{0.56,0.35,0.01}{\textit{#1}}}
\newcommand{\CommentVarTok}[1]{\textcolor[rgb]{0.56,0.35,0.01}{\textbf{\textit{#1}}}}
\newcommand{\ConstantTok}[1]{\textcolor[rgb]{0.00,0.00,0.00}{#1}}
\newcommand{\ControlFlowTok}[1]{\textcolor[rgb]{0.13,0.29,0.53}{\textbf{#1}}}
\newcommand{\DataTypeTok}[1]{\textcolor[rgb]{0.13,0.29,0.53}{#1}}
\newcommand{\DecValTok}[1]{\textcolor[rgb]{0.00,0.00,0.81}{#1}}
\newcommand{\DocumentationTok}[1]{\textcolor[rgb]{0.56,0.35,0.01}{\textbf{\textit{#1}}}}
\newcommand{\ErrorTok}[1]{\textcolor[rgb]{0.64,0.00,0.00}{\textbf{#1}}}
\newcommand{\ExtensionTok}[1]{#1}
\newcommand{\FloatTok}[1]{\textcolor[rgb]{0.00,0.00,0.81}{#1}}
\newcommand{\FunctionTok}[1]{\textcolor[rgb]{0.00,0.00,0.00}{#1}}
\newcommand{\ImportTok}[1]{#1}
\newcommand{\InformationTok}[1]{\textcolor[rgb]{0.56,0.35,0.01}{\textbf{\textit{#1}}}}
\newcommand{\KeywordTok}[1]{\textcolor[rgb]{0.13,0.29,0.53}{\textbf{#1}}}
\newcommand{\NormalTok}[1]{#1}
\newcommand{\OperatorTok}[1]{\textcolor[rgb]{0.81,0.36,0.00}{\textbf{#1}}}
\newcommand{\OtherTok}[1]{\textcolor[rgb]{0.56,0.35,0.01}{#1}}
\newcommand{\PreprocessorTok}[1]{\textcolor[rgb]{0.56,0.35,0.01}{\textit{#1}}}
\newcommand{\RegionMarkerTok}[1]{#1}
\newcommand{\SpecialCharTok}[1]{\textcolor[rgb]{0.00,0.00,0.00}{#1}}
\newcommand{\SpecialStringTok}[1]{\textcolor[rgb]{0.31,0.60,0.02}{#1}}
\newcommand{\StringTok}[1]{\textcolor[rgb]{0.31,0.60,0.02}{#1}}
\newcommand{\VariableTok}[1]{\textcolor[rgb]{0.00,0.00,0.00}{#1}}
\newcommand{\VerbatimStringTok}[1]{\textcolor[rgb]{0.31,0.60,0.02}{#1}}
\newcommand{\WarningTok}[1]{\textcolor[rgb]{0.56,0.35,0.01}{\textbf{\textit{#1}}}}
\usepackage{longtable,booktabs}
% Correct order of tables after \paragraph or \subparagraph
\usepackage{etoolbox}
\makeatletter
\patchcmd\longtable{\par}{\if@noskipsec\mbox{}\fi\par}{}{}
\makeatother
% Allow footnotes in longtable head/foot
\IfFileExists{footnotehyper.sty}{\usepackage{footnotehyper}}{\usepackage{footnote}}
\makesavenoteenv{longtable}
\usepackage{graphicx,grffile}
\makeatletter
\def\maxwidth{\ifdim\Gin@nat@width>\linewidth\linewidth\else\Gin@nat@width\fi}
\def\maxheight{\ifdim\Gin@nat@height>\textheight\textheight\else\Gin@nat@height\fi}
\makeatother
% Scale images if necessary, so that they will not overflow the page
% margins by default, and it is still possible to overwrite the defaults
% using explicit options in \includegraphics[width, height, ...]{}
\setkeys{Gin}{width=\maxwidth,height=\maxheight,keepaspectratio}
% Set default figure placement to htbp
\makeatletter
\def\fps@figure{htbp}
\makeatother
\setlength{\emergencystretch}{3em} % prevent overfull lines
\providecommand{\tightlist}{%
  \setlength{\itemsep}{0pt}\setlength{\parskip}{0pt}}
\setcounter{secnumdepth}{5}
\usepackage{booktabs}
\usepackage[]{natbib}
\bibliographystyle{apalike}

\title{Marketing Research}
\author{Mike Nguyen}
\date{2021-01-05}

\begin{document}
\maketitle

{
\setcounter{tocdepth}{1}
\tableofcontents
}
\hypertarget{prerequisites}{%
\chapter{Prerequisites}\label{prerequisites}}

\hypertarget{intro}{%
\chapter{Introduction}\label{intro}}

You can label chapter and section titles using \texttt{\{\#label\}} after them, e.g., we can reference Chapter \ref{intro}. If you do not manually label them, there will be automatic labels anyway, e.g., Chapter \ref{methods}.

Figures and tables with captions will be placed in \texttt{figure} and \texttt{table} environments, respectively.

\begin{Shaded}
\begin{Highlighting}[]
\KeywordTok{par}\NormalTok{(}\DataTypeTok{mar =} \KeywordTok{c}\NormalTok{(}\DecValTok{4}\NormalTok{, }\DecValTok{4}\NormalTok{, }\FloatTok{.1}\NormalTok{, }\FloatTok{.1}\NormalTok{))}
\KeywordTok{plot}\NormalTok{(pressure, }\DataTypeTok{type =} \StringTok{'b'}\NormalTok{, }\DataTypeTok{pch =} \DecValTok{19}\NormalTok{)}
\end{Highlighting}
\end{Shaded}

\begin{figure}

{\centering \includegraphics[width=0.8\linewidth]{marketing_research_files/figure-latex/nice-fig-1} 

}

\caption{Here is a nice figure!}\label{fig:nice-fig}
\end{figure}

Reference a figure by its code chunk label with the \texttt{fig:} prefix, e.g., see Figure \ref{fig:nice-fig}. Similarly, you can reference tables generated from \texttt{knitr::kable()}, e.g., see Table \ref{tab:nice-tab}.

\begin{Shaded}
\begin{Highlighting}[]
\NormalTok{knitr}\OperatorTok{::}\KeywordTok{kable}\NormalTok{(}
  \KeywordTok{head}\NormalTok{(iris, }\DecValTok{20}\NormalTok{), }\DataTypeTok{caption =} \StringTok{'Here is a nice table!'}\NormalTok{,}
  \DataTypeTok{booktabs =} \OtherTok{TRUE}
\NormalTok{)}
\end{Highlighting}
\end{Shaded}

\begin{table}

\caption{\label{tab:nice-tab}Here is a nice table!}
\centering
\begin{tabular}[t]{rrrrl}
\toprule
Sepal.Length & Sepal.Width & Petal.Length & Petal.Width & Species\\
\midrule
5.1 & 3.5 & 1.4 & 0.2 & setosa\\
4.9 & 3.0 & 1.4 & 0.2 & setosa\\
4.7 & 3.2 & 1.3 & 0.2 & setosa\\
4.6 & 3.1 & 1.5 & 0.2 & setosa\\
5.0 & 3.6 & 1.4 & 0.2 & setosa\\
\addlinespace
5.4 & 3.9 & 1.7 & 0.4 & setosa\\
4.6 & 3.4 & 1.4 & 0.3 & setosa\\
5.0 & 3.4 & 1.5 & 0.2 & setosa\\
4.4 & 2.9 & 1.4 & 0.2 & setosa\\
4.9 & 3.1 & 1.5 & 0.1 & setosa\\
\addlinespace
5.4 & 3.7 & 1.5 & 0.2 & setosa\\
4.8 & 3.4 & 1.6 & 0.2 & setosa\\
4.8 & 3.0 & 1.4 & 0.1 & setosa\\
4.3 & 3.0 & 1.1 & 0.1 & setosa\\
5.8 & 4.0 & 1.2 & 0.2 & setosa\\
\addlinespace
5.7 & 4.4 & 1.5 & 0.4 & setosa\\
5.4 & 3.9 & 1.3 & 0.4 & setosa\\
5.1 & 3.5 & 1.4 & 0.3 & setosa\\
5.7 & 3.8 & 1.7 & 0.3 & setosa\\
5.1 & 3.8 & 1.5 & 0.3 & setosa\\
\bottomrule
\end{tabular}
\end{table}

You can write citations, too. For example, we are using the \textbf{bookdown} package \citep{R-bookdown} in this sample book, which was built on top of R Markdown and \textbf{knitr} \citep{xie2015}.

\hypertarget{advertising}{%
\chapter{Advertising}\label{advertising}}

Prior distribution for TV advertising elasticity for consumer packaged goods can be found in \citep{Shapiro_2018}.
A substantial portion of products has statistically insignificant or negative estimates for advertising elasticity. TV ads might not be the best vehicle to reach the customer now.

\hypertarget{brand-equity}{%
\chapter{Brand Equity}\label{brand-equity}}

\hypertarget{celebrity-endorsement}{%
\chapter{Celebrity Endorsement}\label{celebrity-endorsement}}

Based on \href{https://bookdown.org/mike/social-theory/power-distance-theory.html}{Power Distance Theory}, power distance beliefs moderates the positive effect of celebrity endorsers on evaluations of advertising \citep{Winterich_2018}

\hypertarget{virality}{%
\chapter{Virality}\label{virality}}

The age group from 18-24 can be more challenging to target since they reported fewer positive responses compared to other age groups \citep{Libert_2014}. Since their threshold is higher due to the overwhelming digital content they encounter, their emotions are harder to activate.

Wide gaps between two opposing opinions can increase future discussion (reviews); however, fluctuation around one opinion does not. Higher informational content moderate positively the relationship between opposing opinions and future reviews. \citep{Nagle_2014}

WOM should be taken under the context of the evolution of conversation. \citep{Berger_2014}

\textbf{Valuable virality} \citep{Akpinar_2017}

Establishes causality of appeals and brand integralness to virality.

\begin{itemize}
\tightlist
\item
  Compared to informative appeal ads (focus on product features), emotional appeal ads ( drama, mood, music, emotion-eliciting strategies) are more likely to be shared.\\
\item
  Informative appeals boost brand evaluations and purchase because the brand is an integral part of the ad content.\\
\item
  Hence, to combine the best of both worlds, emotional brand-integral ads boost sharing while also bolstering brand-related outcomes (embedded the brand into the stories).\\
\item
  The mechanism that ads affect brand evaluation through two simultaneous mechanisms: persuasive attempt and brand knowledge.
\end{itemize}

``keeping viewers involved depends in large part on two emotions: joy and surprise''(Teixeira, 2012) - ``Research: The Emotions that Make Marketing Campaigns Go Viral''

``ads that produce stable emotional states generally aren't effective at engaging viewers for very long.'' And ``shock may get people to watch an ad privately, it often works against their desire to share the spot. Like Clothing Drive that try to mimic ``Swear Jar'' by Bud Light'' (Teixeira, 2012)

(Westbrook, 1987, p.~261) defined WOM as ``informational communication directed at other consumers about the ownership, usage, or characteristics of particular goods and services or their sellers.'' WOM evolves to ``word of mouse'' (or e-WOM).

\citep{Berger_2014}:

WOM is goal-driven and serves five key functions:

\begin{itemize}
\tightlist
\item
  \protect\hyperlink{impression-management}{Impression Management}\\
\item
  \protect\hyperlink{emotion-regulation}{Emotion regulation}\\
\item
  \protect\hyperlink{information-acquisition}{Information acquisition}\\
\item
  \protect\hyperlink{social-bonding}{Social bonding}\\
\item
  \protect\hyperlink{persuading-others}{Persuasion}
\end{itemize}

The motivation behind this process is self -serving and even subconscious level. Hence, we can predict the types of news and information people are most likely to discuss.

\begin{itemize}
\tightlist
\item
  Over 70\% of everyday speech is about the self, such as personal experiences or relationships (Dunbar et al., 1997). Consistent with previous research, online behaviors exhibit the same pattern: over 70\% of social media posts are about self (Naaman et al., 2010).
\item
  At any given point in time, multiple drivers can be present.
\end{itemize}

\hypertarget{impression-management}{%
\section{Impression Management}\label{impression-management}}

Consumers buy to signal their derived identities to achieve desired impressions (J. Berger \& Heath, 2007)

Interpersonal communication (WOM) influences impression management in 3 ways: (1) self-enhancement, (2) identity -signaling, and (3) filling conversational space.

\begin{itemize}
\tightlist
\item
  \textbf{Self-enhancement}: human has a tendency to self-enhance (Markus, 1987). Hence, people like to share things to help them look good (Chung \& Darke, 2006; Hennig-Thurau et al., 2004). Since bragging too much about oneself might have an opposite effect. People sometimes engage in'' humblebragging'': brag but self-deprecating at the same time. (Sezer et al., 2018)\\
\item
  \textbf{Identity-signalling}: they talk about themselves to signal they have certain traits or expertise (G. Packard \& Wooten, 2013), such as opinion leaders talk to signal their identity (share to show their knowledge).\\
\item
  \textbf{Filling conversational space}: (small talk) people engage in small talk to avoid pauses and silence in between conversations so that the other party would not make lousy inferences about the person. (J. Berger, 2014)
\end{itemize}

Hence impression management induces people to share things that are

\begin{itemize}
\item
  \textbf{Entertaining} (i.e., interesting, surprising, funny or extreme) so that they shared look more interesting, and in-the-know. Interesting products get more word of mouth (J. Berger \& Iyengar, 2013; J. Berger \& Milkman, 2012; J. Berger \& Schwartz, 2011). Moderate controversy is a catalyst for word of mouth to spread (Chen \& Berger, 2013). And sometimes, people even exaggerate stories to make things more interesting (C. Heath, 1996); around 60\% of the stories are distorted (Marsh \& Tversky, 2004). Things that can be considered interesting are those that are either novel, exciting, or violate previous expectations (i.e., surprise) (Berlyne, 2006; Silvia, 2008). Nevertheless, comprehensibility can be a moderate of whether a novel thing is interesting: novel and comprehensible equals interesting, while novel and incomprehensible equal confusing (Silvia, 2008). Novelty refers to things that are new, surprising, exciting, or complex \citep{Berlyne}. Comprehensibility refers to the fact that the novelty must be understandable.
\item
  \textbf{Useful}: because it makes the sharer smart and helpful. Empirical evidence show that useful stories(J. Berger \& Milkman, 2012) and higher quality brands (Lovett et al., 2013) are more likely to be shared.\\
\item
  \textbf{Self-concept relevant}: different people see different domain as the optimal signal amplifier such as politics, economics, etc. hence symbolic products are more likely to be shared than utilitarian ones (Chung \& Darke, 2006) (even though useful information is helpful, but the signal to be smart and helpful are not appreciated as being interesting and entertaining).\\
\item
  \textbf{Status related}: premium brands are talked about more (Lovett et al., 2013). Knowledge is can also be considered as a status symbol, and people share to signal that they re in the know (Ritson \& Elliott, 1999)\\
\item
  \textbf{Unique}: unique products are more likely to be shared, but people with a high need for uniqueness will be less likely to share positive WOM because they do not want people to adopt the product and reduces their uniqueness (Cheema \& Kaikati, 2010). Sometimes they even scare others by citing product complexity (Moldovan et al., 2015)\\
\item
  \textbf{Common ground}: sharing common things induces listener to convey that there is interpersonal similarity and facilitates conversation.\\
\item
  \textbf{Emotional valence}: people prefer to share positive things than negative news (J. Berger \& Milkman, 2012). However, in some special cases, negative WOM can be perceived as a desired component such as reviewers are seen as more intelligent, (Amabile, 1983). Which is moderated by whether people are refereeing to themselves or others will affect WOM valence (positive vs.~negative)(Kamins et al., 1997). Positive WOM when talking oneself conveys a positive image, while negative WOM about others convey that they are relatively better than the other party.\\
\item
  \textbf{Incidental arousal}: unrelated arousal (e.g., running in place) can increase sharing in general (J. Berger, 2011) because of the increase in arousal levels.\\
\item
  \textbf{Accessible things}: products with more cued or triggered will be discussed more (J. Berger \& Schwartz, 2011). Hence, more advertised products generate more WOM (Onishi \& Manchanda, 2012). ``top of mind'' and ``tip of the tongue.'' Availability bias means that the easier we can recall something, the more likely we are going to talk bout it. Publicly visible products also have more WOM (J. Berger \& Schwartz, 2011)
\end{itemize}

\hypertarget{emotion-regulation}{%
\section{Emotion regulation}\label{emotion-regulation}}

WOM help consumers regulate their emotions. Emotion regulation refers to ``a person's ability to effectively manage and respond to an emotional experience'' (Rolston \& Lloyd-Richardson, 2017). and sharing with other (WOM) can facilitate emotion regulation by:

\begin{itemize}
\tightlist
\item
  \textbf{Generating social support}: sharing can give you comfort and consolation (Rimé, 2009). For example, after a negative emotional experience, sharing can help increase well-being through perceived social support (J. A. Berger \& Buechel, 2012)\\
\item
  \textbf{Venting}: WOM (sharing) helps people achieve catharsis when experiencing negative experiences(Alicke et al., 1992): consumers express when they are angry (Wetzer et al., 2007) or dissatisfied (Anderson, 1998).\\
\item
  \textbf{Sensemaking}: ``talking can help people understand what they feel and why'' (J. Berger, 2014)\\
\item
  \textbf{Reducing dissonance}: even after making a decision, we want to make sure we make the right choice by talking to others tot reduce feelings of doubt (Rosnow, 1980)\\
\item
  \textbf{Taking vengeance:} to punish the company (different from venting, which makes you feel better), consumers share negative consumption experience (Ward \& Ostrom, 2006)\\
\item
  \textbf{Rehearsal}: people talk to relive positive experiences (Rimé, 2009)
\end{itemize}

Emotion regulation drives people to share (1) emotional content (2) influence the valence of the content shared, (3) lead people to share more emotionally arousing content:

\begin{itemize}
\tightlist
\item
  \textbf{Emotionality}: more emotional intensity, more sharing (J. Berger \& Milkman, 2012), people share more emotional social anecdotes (Peters et al., 2009). But not all emotions increase sharing: for example, shame and guilt decreases sharing (Finkenauer \& Rimé, 1998) since sharing those things makes people look bad.\\
\item
  \textbf{Valence}: ``Emotion regulation tends to focus on the management of negative emotions'' (J. Berger, 2014), but a counter vailing effect is that under impression management, people avoid sharing negative stories since it might reflect on the shares. Sharing negative things can decrease willingness to share (Chen \& Berger, 2013)\\
\item
  \textbf{Emotional arousal}:

  \begin{itemize}
  \tightlist
  \item
    \textbf{Negative} side: low arousal emotion (sadness), high arousal (anger or anxiety) should increase the need to vent.\\
  \item
    \textbf{Positive} side: low arousal (contentment), high arousal (awe, excitement or amusement) increase desires for rehearsal.
  \end{itemize}
\end{itemize}

\hypertarget{information-acquisition}{%
\section{Information acquisition}\label{information-acquisition}}

WOM helps acquire information.

\begin{itemize}
\tightlist
\item
  \textbf{Seeking advice}: (Hennig-Thurau et al., 2004; Rimé, 2009). Rather than trial and error, or direct observation, gossip serves as another form of learning (Baumeister et al., 2004)\\
\item
  \textbf{Resolving problems}: online forums or other online opinion platforms help customers find a solution quickly from others (Hennig-Thurau et al., 2004)
\end{itemize}

Hence the process of information acquisition would induce to talk more about

\begin{itemize}
\tightlist
\item
  \textbf{Risky, important, complex, or uncertainty-ridden decision}: talking to others can reduce risk, simplify complexity, and increase consumer's confidence (Hennig-Thurau et al., 2004)\\
\item
  \textbf{Lack of (trustworthy) information}:
\end{itemize}

\hypertarget{social-bonding}{%
\section{Social bonding}\label{social-bonding}}

``people have a fundamental desire for social relationships (Baumeister \& Leary, 1995), and interpersonal communication fills that need (Hennig-Thurau et al., 2004)''. Empirically, people in brand communities because want to connect with others like them (Muniz \& O'Guinn, 2001). Sharing facilitates social bonding through

\begin{itemize}
\tightlist
\item
  \textbf{Reinforce shared views}: group memberships, and one's place in the social hierarchy. WOM (sharing) allows people to fit in as a group member (J. Berger \& Heath, 2007; Ritson \& Elliott, 1999)\\
\item
  \textbf{Reducing loneliness and social exclusion}: sharing decrease interpersonal distance and decreases loneliness and exclusion (Maner et al., 2007)
\end{itemize}

Social bonding drives what people share: people share things are :

\begin{itemize}
\tightlist
\item
  \textbf{Common ground}: social bonding drives people to talk about things they have in common because it makes them feel socially connected (Clark \& Kashima, 2007)
\item
  \textbf{Emotionality}: ``sharing an emotional story or narrative increases the chance that others will feel similarly'' (J. Berger, 2014). Emotional similarity increase group cohesiveness (Barsade \& Gibson, 2007). And emotion sharing and social bonding can be a simultaneous process since previous works found that emotion sharing bonds people (Peters \& Kashima, 2007) and high arousal emotion increase social bonding needs (Chen \& Berger, 2013)
\end{itemize}

\hypertarget{persuading-others}{%
\section{Persuading others}\label{persuading-others}}

people use WOM to persuade others to affect their satisfaction or choices (J. Berger, 2014). Persuasive drive people to share things that are

\begin{itemize}
\tightlist
\item
  \textbf{Emotionally polarized} (polarized valence): since the goal is to convince something is good or bad (i.e., people share extremely rather than moderately positive(negative) information\\
\item
  \textbf{Arousing content}: ``people who want to persuade others may find shared arousing content to incite others to take desired actions'' (J. Berger, 2014)
\end{itemize}

\hypertarget{moderators}{%
\section{Moderators}\label{moderators}}

\begin{itemize}
\tightlist
\item
  Two key moderators when different WOM functions have a greater impact:

  \begin{itemize}
  \tightlist
  \item
    \textbf{Communication audience}:

    \begin{itemize}
    \tightlist
    \item
      Tie strength
    \item
      Audience size
    \item
      Tie status
    \end{itemize}
  \item
    \textbf{Communication Channel}

    \begin{itemize}
    \tightlist
    \item
      Written vs.~oral
    \item
      Identifiability: whether communicators are identifiable. (anonymously)\\
    \item
      Audience salience: whether the audience is salient during communication (online communication, the audience is less salient).
    \end{itemize}
  \end{itemize}
\end{itemize}

Social motivation: why we share :\citep{why_some_videos_go_viral_2015}

\begin{itemize}
\tightlist
\item
  Impression management:

  \begin{itemize}
  \tightlist
  \item
    authority (e.g., ``I want to demonstrate my knowledge'')\\
  \item
    Social utility (e.g., this could be useful to my friends)\\
  \item
    coolhunter (e.g., I want to be the first o tell my friends)\\
  \item
    Zeitgeist (e.g., It's about a current trend or event)\\
  \item
    Conversation starting (e.g., I want to start an online conversation)\\
  \item
    Self-expression (e.g., it says something about me)\\
  \item
    Social good (e.g., It's for a good cause, and I want to help'').\\
  \end{itemize}
\item
  Information acquisition:

  \begin{itemize}
  \tightlist
  \item
    Opinion seeking (e.g., ``I want to see what my friends think)\\
  \end{itemize}
\item
  Social bonding:

  \begin{itemize}
  \tightlist
  \item
    Shared passion (e.g, ``it lets me connect with my friends bout a shared interest)\\
  \item
    Social in real life (e.g., it will help me socialize with my friends offline).
  \end{itemize}
\end{itemize}

More frequent exposure to perceptually or conceptually related cues increases product accessibility and makes the product easier to process. In turn, this increased accessibility influences product evaluation and choice, which are found to vary directly with the frequency of exposure to conceptually related cues. These results support the hypothesis that conceptual priming effects can have a strong impact on real-world consumer judgment'' \citep{Berger_2008}

\hypertarget{finance-marketing-interface}{%
\chapter{Finance-Marketing Interface}\label{finance-marketing-interface}}

\hypertarget{tobins-q}{%
\section{Tobin's q}\label{tobins-q}}

Retrieve data from WRDS

\begin{Shaded}
\begin{Highlighting}[]
\CommentTok{# to set up connection from R to WRDS (https://wrds-www.wharton.upenn.edu/pages/support/programming-wrds/programming-r/r-from-your-computer/)}
\KeywordTok{library}\NormalTok{(RPostgres)}
\KeywordTok{library}\NormalTok{(dplyr)}

\CommentTok{# I've set up wrds connection before hand. Please use your username and password here.}

\CommentTok{# wrds <- dbConnect(Postgres(),}
\CommentTok{#                   host='wrds-pgdata.wharton.upenn.edu',}
\CommentTok{#                   port=9737,}
\CommentTok{#                   dbname='wrds',}
\CommentTok{#                   sslmode='require',}
\CommentTok{#                   user='')}

\CommentTok{# Check variables (column headers) in COMP ANNUAL FUNDAMENTAL}
\CommentTok{#uses the already-established wrds connection to prepare the SQL query string and save the query as the result res.}
\CommentTok{# check avaiable databases: https://wrds-web.wharton.upenn.edu/wrds/tools/variable.cfm?vendor_id=7}
\NormalTok{res <-}\StringTok{ }\KeywordTok{dbSendQuery}\NormalTok{(wrds, }\StringTok{"select column_name}
\StringTok{                   from information_schema.columns}
\StringTok{                   where table_schema='compa'}
\StringTok{                   and table_name='funda'}
\StringTok{                   order by column_name"}\NormalTok{)}
\NormalTok{data <-}\StringTok{ }\KeywordTok{dbFetch}\NormalTok{(res, }\DataTypeTok{n=}\OperatorTok{-}\DecValTok{1}\NormalTok{) }\CommentTok{# fetches the data that results from running the query res against wrds and stores it as data}
\KeywordTok{dbClearResult}\NormalTok{(res) }\CommentTok{# closes the connection}
\KeywordTok{head}\NormalTok{(data)}
\end{Highlighting}
\end{Shaded}

\begin{verbatim}
##   column_name
## 1       acchg
## 2        acco
## 3       accrt
## 4     acctchg
## 5     acctstd
## 6        acdo
\end{verbatim}

\begin{Shaded}
\begin{Highlighting}[]
\CommentTok{# select everything}
\NormalTok{res <-}\StringTok{ }\KeywordTok{dbSendQuery}\NormalTok{(wrds, }\StringTok{"select * from compa.funda"}\NormalTok{)}

\CommentTok{# from compa.funda}

\CommentTok{# only select the following variables}
\NormalTok{res <-}\StringTok{ }\KeywordTok{dbSendQuery}\NormalTok{(wrds, }\StringTok{"select gvkey, datadate, fyear, indfmt, consol, popsrc, datafmt, tic, cusip, conm, curcd, fyr, act, at, bkvlps, ceq, ch, che, dltt, dlc, emp, np, exchg, cik, costat, naicsh,mkvalt from compa.funda"}\NormalTok{) }\CommentTok{#check variables from (https://wrds-web.wharton.upenn.edu/wrds/ds/comp/funda/index.cfm?navId=80)}
\end{Highlighting}
\end{Shaded}

\begin{verbatim}
## Warning in result_create(conn@ptr, statement): Closing open result set,
## cancelling previous query
\end{verbatim}

\begin{Shaded}
\begin{Highlighting}[]
\NormalTok{data1 <-}\StringTok{ }\KeywordTok{dbFetch}\NormalTok{(res, }\DataTypeTok{n=}\OperatorTok{-}\DecValTok{1}\NormalTok{)}
\KeywordTok{dbClearResult}\NormalTok{(res)}

\NormalTok{data =}\StringTok{ }\NormalTok{data1 }\OperatorTok{%>%}
\StringTok{    }\KeywordTok{distinct}\NormalTok{(gvkey,datadate,fyear,tic,conm,}\DataTypeTok{.keep_all =}\NormalTok{ T)}
\end{Highlighting}
\end{Shaded}

Calculate Tobin's Q

\begin{Shaded}
\begin{Highlighting}[]
\NormalTok{tobin_q_dt =}\StringTok{ }\NormalTok{data }\OperatorTok{%>%}
\StringTok{    }\KeywordTok{mutate}\NormalTok{(}\DataTypeTok{mkvalt =} \KeywordTok{coalesce}\NormalTok{(mkvalt,}\DecValTok{0}\NormalTok{),}
           \DataTypeTok{dltt =} \KeywordTok{coalesce}\NormalTok{(dltt,}\DecValTok{0}\NormalTok{),}
           \DataTypeTok{at =} \KeywordTok{coalesce}\NormalTok{(at,}\DecValTok{0}\NormalTok{),}
           \DataTypeTok{dlc =} \KeywordTok{coalesce}\NormalTok{(dlc,}\DecValTok{0}\NormalTok{),}
           \DataTypeTok{act =} \KeywordTok{coalesce}\NormalTok{(act,}\DecValTok{0}\NormalTok{)) }\OperatorTok{%>%}
\StringTok{    }\KeywordTok{mutate}\NormalTok{(}\DataTypeTok{tobin_q =}\NormalTok{ (mkvalt }\OperatorTok{+}\StringTok{ }\KeywordTok{ifelse}\NormalTok{((dlc }\OperatorTok{-}\StringTok{ }\NormalTok{act) }\OperatorTok{>=}\DecValTok{0}\NormalTok{,}\KeywordTok{as.numeric}\NormalTok{(dlc}\OperatorTok{-}\NormalTok{act),}\DecValTok{0}\NormalTok{) }\OperatorTok{+}\StringTok{ }\NormalTok{dltt)}\OperatorTok{/}\NormalTok{at ) }\OperatorTok{%>%}\StringTok{  }\CommentTok{#follow Chung (1994) (sum of stocks (preferred + common) + debt(short-term liabilities - short-term assets + long-term debt))/(total assets) # note: take only excess of short-term liabilities over short-term assets to be included in debt. }
\StringTok{    }\KeywordTok{select}\NormalTok{(tobin_q,gvkey,datadate,fyear,conm)}
\KeywordTok{head}\NormalTok{(tobin_q_dt)}
\end{Highlighting}
\end{Shaded}

\begin{verbatim}
##     tobin_q  gvkey   datadate fyear                  conm
## 1       Inf 001000 1961-12-31  1961 A & E PLASTIK PAK INC
## 2       NaN 001000 1962-12-31  1962 A & E PLASTIK PAK INC
## 3       Inf 001000 1963-12-31  1963 A & E PLASTIK PAK INC
## 4 0.3686441 001000 1964-12-31  1964 A & E PLASTIK PAK INC
## 5 0.4995671 001000 1965-12-31  1965 A & E PLASTIK PAK INC
## 6 0.4563786 001000 1966-12-31  1966 A & E PLASTIK PAK INC
\end{verbatim}

\hypertarget{event-study}{%
\section{Event study}\label{event-study}}

\hypertarget{abnormal-return}{%
\subsection{Abnormal Return}\label{abnormal-return}}

\url{https://rpubs.com/gregoryarobinson/315096}

\url{https://www.eventstudytools.com/introduction-event-study-methodology}

\begin{Shaded}
\begin{Highlighting}[]
\KeywordTok{library}\NormalTok{(EventStudy)}
\end{Highlighting}
\end{Shaded}

Another package

\url{https://cran.r-project.org/web/packages/eventstudies/vignettes/replications.pdf}

\begin{Shaded}
\begin{Highlighting}[]
\KeywordTok{library}\NormalTok{(eventstudies)}
\end{Highlighting}
\end{Shaded}

Replication of \citep{Chen_2004}

\begin{Shaded}
\begin{Highlighting}[]
\KeywordTok{data}\NormalTok{(TerrorIndiceReturns)}
\KeywordTok{data}\NormalTok{(TerrorAttack)}
\NormalTok{TerrorIndiceCAR <-}\StringTok{ }\KeywordTok{lapply}\NormalTok{(}\DecValTok{1}\OperatorTok{:}\StringTok{ }\KeywordTok{ncol}\NormalTok{(TerrorIndiceReturns), }\ControlFlowTok{function}\NormalTok{(x)\{}
\CommentTok{# 10-day window around the event}
\NormalTok{event <-}\StringTok{ }\KeywordTok{phys2eventtime}\NormalTok{(}\KeywordTok{na.omit}\NormalTok{(TerrorIndiceReturns[ , x,}
\DataTypeTok{drop =} \OtherTok{FALSE}\NormalTok{]),}
\NormalTok{TerrorAttack,}
\DecValTok{10}\NormalTok{)}
\CommentTok{# Estimate ARs}
\NormalTok{esMean <-}\StringTok{ }\KeywordTok{constantMeanReturn}\NormalTok{(event}\OperatorTok{$}\NormalTok{z.e[}\KeywordTok{which}\NormalTok{(}\KeywordTok{attributes}\NormalTok{(event}\OperatorTok{$}\NormalTok{z.e)}\OperatorTok{$}\NormalTok{index}
\OperatorTok{%in%}\StringTok{ }\DecValTok{-30}\OperatorTok{:-}\DecValTok{11}\NormalTok{), ],}
\DataTypeTok{residual =} \OtherTok{FALSE}\NormalTok{)}
\NormalTok{ar <-}\StringTok{ }\NormalTok{event}\OperatorTok{$}\NormalTok{z.e }\OperatorTok{-}\StringTok{ }\NormalTok{esMean}
\NormalTok{ar <-}\StringTok{ }\KeywordTok{window}\NormalTok{(ar, }\DataTypeTok{start =} \DecValTok{0}\NormalTok{, }\DataTypeTok{end =} \DecValTok{10}\NormalTok{)}
\CommentTok{# CAR}
\NormalTok{car <-}\StringTok{ }\KeywordTok{remap.cumsum}\NormalTok{(ar, }\DataTypeTok{base =} \KeywordTok{as.numeric}\NormalTok{(ar[}\DecValTok{1}\NormalTok{, }\DecValTok{1}\NormalTok{]))}
\KeywordTok{names}\NormalTok{(car) <-}\StringTok{ }\KeywordTok{colnames}\NormalTok{(TerrorIndiceReturns[ , x,}
\DataTypeTok{drop =} \OtherTok{FALSE}\NormalTok{])}
\KeywordTok{return}\NormalTok{(car)}
\NormalTok{\})}
\KeywordTok{names}\NormalTok{(TerrorIndiceCAR) <-}\StringTok{ }\KeywordTok{colnames}\NormalTok{(TerrorIndiceReturns)}
\CommentTok{# Compile for all indices}
\NormalTok{TerrorIndiceCAR <-}\StringTok{ }\KeywordTok{do.call}\NormalTok{(cbind, TerrorIndiceCAR)}
\CommentTok{# 11-day CAR}
\NormalTok{TerrorIndiceCAR[}\DecValTok{11}\NormalTok{, ]}
\end{Highlighting}
\end{Shaded}

\begin{verbatim}
##       S.P500       Dow      NYSE    NASDAQ   Toronto   Mexico    London
## 10 -3.832794 -7.897366 -3.981837 -9.993519 -9.110412 -6.53973 -8.602633
##      Germany Eurostoxx50    France    Spain Switzerland   Austria     Italy
## 10 -10.45718   -7.596024 -9.518721 -10.6724   -7.292201 -7.894399 -15.21462
##      Belgium Amsterdam Portugal  Helsinki    Norway    Sweden     Tokyo
## 10 -9.219016 -10.82756 1.698079 -5.500048 -12.39281 -4.691903 -3.603373
##     HongKong   S.Korea     India   Jakarta Singapore Kualalampur Australia
## 10 -5.231521 -16.64571 -17.02777 -9.933654 -17.98219   -14.55417 -8.597988
##          NZL  Pakistan     Saudi   Israel     Jberg
## 10 -5.635766 -12.42189 -11.32466 -7.81159 -12.59022
\end{verbatim}

\begin{Shaded}
\begin{Highlighting}[]
\NormalTok{TerrorIndiceCAR[}\DecValTok{6}\NormalTok{, ] }\CommentTok{# 6-day CAR}
\end{Highlighting}
\end{Shaded}

\begin{verbatim}
##     S.P500       Dow      NYSE    NASDAQ  Toronto    Mexico    London   Germany
## 5 -7.72379 -10.56695 -8.094739 -10.13781 -7.01292 -13.17444 -4.533527 -8.063819
##   Eurostoxx50    France     Spain Switzerland   Austria     Italy   Belgium
## 5   -5.940704 -8.707865 -9.426443   -5.967055 -4.428614 -13.98185 -8.510382
##   Amsterdam  Portugal  Helsinki    Norway   Sweden      Tokyo  HongKong
## 5 -8.521101 -6.875878 -4.802975 -9.894391 -4.96208 -0.5636291 -5.570385
##     S.Korea     India  Jakarta Singapore Kualalampur Australia       NZL
## 5 -11.81901 -11.77841 -4.91951 -13.74157   -10.50214 -6.814062 -7.101218
##   Pakistan     Saudi    Israel     Jberg
## 5 -10.3741 -5.320904 -11.03399 -12.09548
\end{verbatim}

\hypertarget{advertising-1}{%
\chapter{Advertising}\label{advertising-1}}

People don't like intrusive ads, especially when they are vulnerable \citep[p.143]{kellogg_2005}

  \bibliography{book.bib,packages.bib,references.bib}

\end{document}
